% GENERATED FILE. Edit canonical lessons, then run npm run generate:books.
% canonical-source-hash: fnv1a64:961f51a3948e1d00
% canonical-lessons: MR-C175-vicitra, MR-C175-vegla, MR-C175-sarkha, MR-C175-yogya, MR-C175-ayogya

\chapter{Strange, Different, Similar, Proper, Wrong}
\label{ch:mr-strange-different-similar-proper-wrong}
\hlchaptermodality{175}

\begin{chapteropening}
\textbf{By the end of this chapter:} \emph{I can say strange, different, similar, proper and wrong.}

Complete the last lesson of chapter 175: I can say strange, different, similar, proper and wrong.
\end{chapteropening}

\section[\texorpdfstring{\mr{विचित्र} (vicitra)}{विचित्र (vicitra)}]{\mr{विचित्र} (vicitra) — strange}

\label{lesson:MR-C175-vicitra}



\begin{quote}
Before the new one: say the Marathi for unpleasant, then the Marathi for enthusiastic.
\end{quote}

\subsection*{You'll want to know: \mr{विचित्र}}
\textbf{\mr{विचित्र}} — \emph{vicitra} — \textquotedblleft{}strange\textquotedblright{}.

\begin{grammarlens}[title={What you've built}]
One more word for this chapter.
\end{grammarlens}

\subsection*{Your turn}
\begin{itemize}
\raggedright
  \item \emph{Say it:} \emph{vicitra}
  \item \emph{Say it:} \emph{vicitra}, once more
  \item \emph{Say it:} say \emph{vicitra}
  \item \emph{Recall:} say the Marathi for unpleasant, then the Marathi for enthusiastic, then say \emph{vicitra} again
\end{itemize}

\begin{tcolorbox}[breakable,colback=teal!4,colframe=teal!35!black,title={Before you move on}]
What does \mr{विचित्र} mean? (\textquotedblleft{}Strange\textquotedblright{}.) Say it once more.
\end{tcolorbox}
\section[\texorpdfstring{\mr{वेगळा} / \mr{वेगळी} (vegḷā / vegḷī)}{वेगळा / वेगळी (vegḷā / vegḷī)}]{\mr{वेगळा} / \mr{वेगळी} (vegḷā / vegḷī) — different}

\label{lesson:MR-C175-vegla}



\begin{quote}
Before the new one: say the Marathi for enthusiastic, then the Marathi for strange.
\end{quote}

\subsection*{You'll want to know: \mr{वेगळा} / \mr{वेगळी}}
\textbf{\mr{वेगळा} / \mr{वेगळी}} — \emph{vegḷā / vegḷī} — \textquotedblleft{}different\textquotedblright{}.

\begin{grammarlens}[title={What you've built}]
One more word for this chapter.
\end{grammarlens}

\subsection*{Your turn}
\begin{itemize}
\raggedright
  \item \emph{Say it:} \emph{vegḷā / vegḷī}
  \item \emph{Say it:} \emph{vegḷā / vegḷī}, once more
  \item \emph{Say it:} say \emph{vegḷā / vegḷī}
  \item \emph{Recall:} say the Marathi for enthusiastic, then the Marathi for strange, then say \emph{vegḷā / vegḷī} again
\end{itemize}

\begin{tcolorbox}[breakable,colback=teal!4,colframe=teal!35!black,title={Before you move on}]
What does \mr{वेगळा} / \mr{वेगळी} mean? (\textquotedblleft{}Different\textquotedblright{}.) Say it once more.
\end{tcolorbox}
\section[\texorpdfstring{\mr{सारखा} / \mr{सारखी} (sārkhā / sārkhī)}{सारखा / सारखी (sārkhā / sārkhī)}]{\mr{सारखा} / \mr{सारखी} (sārkhā / sārkhī) — similar, alike}

\label{lesson:MR-C175-sarkha}



\begin{quote}
Before the new one: say the Marathi for strange, then the Marathi for different.
\end{quote}

\subsection*{You'll want to know: \mr{सारखा} / \mr{सारखी}}
\textbf{\mr{सारखा} / \mr{सारखी}} — \emph{sārkhā / sārkhī} — \textquotedblleft{}similar, alike\textquotedblright{}.

\begin{grammarlens}[title={What you've built}]
One more word for this chapter.
\end{grammarlens}

\subsection*{Your turn}
\begin{itemize}
\raggedright
  \item \emph{Say it:} \emph{sārkhā / sārkhī}
  \item \emph{Say it:} \emph{sārkhā / sārkhī}, once more
  \item \emph{Say it:} say \emph{sārkhā / sārkhī}
  \item \emph{Recall:} say the Marathi for strange, then the Marathi for different, then say \emph{sārkhā / sārkhī} again
\end{itemize}

\begin{tcolorbox}[breakable,colback=teal!4,colframe=teal!35!black,title={Before you move on}]
What does \mr{सारखा} / \mr{सारखी} mean? (\textquotedblleft{}Similar, alike\textquotedblright{}.) Say it once more.
\end{tcolorbox}
\section[\texorpdfstring{\mr{योग्य} (yogya)}{योग्य (yogya)}]{\mr{योग्य} (yogya) — proper, right}

\label{lesson:MR-C175-yogya}



\begin{quote}
Before the new one: say the Marathi for different, then the Marathi for similar.
\end{quote}

\subsection*{You'll want to know: \mr{योग्य}}
\textbf{\mr{योग्य}} — \emph{yogya} — \textquotedblleft{}proper, right\textquotedblright{}.

\begin{grammarlens}[title={What you've built}]
One more word for this chapter.
\end{grammarlens}

\subsection*{Your turn}
\begin{itemize}
\raggedright
  \item \emph{Say it:} \emph{yogya}
  \item \emph{Say it:} \emph{yogya}, once more
  \item \emph{Say it:} say \emph{yogya}
  \item \emph{Recall:} say the Marathi for different, then the Marathi for similar, then say \emph{yogya} again
\end{itemize}

\begin{tcolorbox}[breakable,colback=teal!4,colframe=teal!35!black,title={Before you move on}]
What does \mr{योग्य} mean? (\textquotedblleft{}Proper, right\textquotedblright{}.) Say it once more.
\end{tcolorbox}
\section[\texorpdfstring{\mr{अयोग्य} (ayogya)}{अयोग्य (ayogya)}]{\mr{अयोग्य} (ayogya) — wrong, improper}

\label{lesson:MR-C175-ayogya}



\begin{quote}
Before the new one: say the Marathi for similar, then the Marathi for proper.
\end{quote}

\subsection*{You'll want to know: \mr{अयोग्य}}
\textbf{\mr{अयोग्य}} — \emph{ayogya} — \textquotedblleft{}wrong, improper\textquotedblright{}.

\begin{grammarlens}[title={What you've built}]
One more word for this chapter.
\end{grammarlens}

\subsection*{Your turn}
\begin{itemize}
\raggedright
  \item \emph{Say it:} \emph{ayogya}
  \item \emph{Say it:} \emph{ayogya}, once more
  \item \emph{Say it:} say \emph{ayogya}
  \item \emph{Recall:} say the Marathi for similar, then the Marathi for proper, then say \emph{ayogya} again
\end{itemize}

\begin{tcolorbox}[breakable,colback=teal!4,colframe=teal!35!black,title={Before you move on}]
What does \mr{अयोग्य} mean? (\textquotedblleft{}Wrong, improper\textquotedblright{}.) Say it once more.
\end{tcolorbox}
